\clearpage
\section{Methodology}
\label{sec:methodology}

This chapter formalises shared--private multi-omics learning, describes the MOCSS and OMIDIENT baselines, introduces Methods~1--3, and defines the evaluation protocol with emphasis on~$\Delta\mathrm{X}\,R^{2}$.

\subsection{Problem formalisation}

Let $M=3$ matched omics modalities be observed for $N$ samples. For modality $m\in\{1,2,3\}$, write $X_m\in\mathbb{R}^{N\times d_m}$. A shared--private encoder produces
\begin{equation}
(S_m,P_m)=E_m(X_m),\qquad S_m,P_m\in\mathbb{R}^{N\times d},
\end{equation}
with embedding dimension $d=32$ in all experiments. Intended semantics are:
\begin{itemize}
\item $S_m$ should capture information about sample~$i$ that is predictable across modalities (cross-view biology, including subtype-related programmes);
\item $P_m$ should capture residual view-specific variation useful for reconstructing $X_m$ but not for explaining other views.
\end{itemize}
Downstream unsupervised subtyping is performed by $k$-means on $\mathrm{mean}_m(S_m)$, on $\mathrm{concat}_m(P_m)$, and on $\mathrm{concat}(P_1,P_2,P_3,\mathrm{mean}_m(S_m))$, among other probes. The scientific claim of interest is not only that some embedding clusters by subtype, but that~$S$ and~$P$ realise the decomposition above.

\paragraph{Shared implementation conventions.}
Unless a method states otherwise, encoders are four-layer multilayer perceptrons with $\tanh$ activations and widths $[512,256,128,32]$. Decoders and predictors use four-layer MLPs with $\tanh$ hidden units and a final sigmoid. Shared codes are passed through a two-layer $\tanh$ projection before contrastive learning. Weights are Kaiming-initialised with zero biases. Reconstruction losses follow the MOCSS-style centred, row-normalised matrix-norm metric rather than a literal mean squared error, ensuring fair comparison with the baselines. The contrastive term is the sum of pairwise InfoNCE losses among the three projected shared embeddings in a batch~\citep{chen2020simple,chen2023mocss}.

\subsection{Baseline: MOCSS}

MOCSS~\citep{chen2023mocss} trains, for each modality, dual reconstruction paths
\begin{equation}
\hat{X}_m^{(S)}=\mathrm{Dec}_{S}(S_m),\qquad
\hat{X}_m^{(P)}=\mathrm{Dec}_{P}(P_m),
\end{equation}
together with a weak elementwise orthogonality penalty between~$\ell_2$-normalised $S_m$ and $P_m$, and InfoNCE on $\mathrm{MLP}(S_m)$ across modalities. Dual full reconstruction is the central inductive bias under scrutiny: each code is encouraged to be individually sufficient for~$X_m$, which conflicts with a residual interpretation of~$P_m$ and invites private codes to store shared subtype signal.

\subsection{Baseline: OMIDIENT}

OMIDIENT~\citep{Safinianaini_Välimäki_Bresson_Gorbonos_Rajamäki_Aaltonen_Marttinen_2025} keeps a MOCSS-style shared autoencoder and contrastive head, but replaces the private pathway with a product-Dirichlet variational encoder--decoder and a corresponding KL regulariser. The hope is that compositional private latents better absorb view-specific residual structure. The shared branch, however, remains essentially MOCSS. This thesis therefore treats OMIDIENT as a private-branch intervention whose shared--private claims must still be verified with the same probes used for AE-private models. A Method~2 hybrid that grafts Dirichlet private codes onto Method~2's joint reconstruction objective is included as a diagnostic companion.

\subsection{Method 1: predictability surrogates}

Method~1 replaces dual solo reconstruction as the definition of shared versus private with a set of predictability constraints. For target modality~$A$ and the other two modalities $\{B,C\}$, five networks are trained:
\begin{align}
h_A(P_A,S_A)&\rightarrow X_A && \text{(within-view joint reconstruction)},\\
f_A(S_B,S_C)&\rightarrow X_A && \text{(cross-shared sufficiency)},\\
g_A(P_B,P_C)&\rightarrow X_A && \text{(cross-private adversary)},\\
k_A(S_A)&\rightarrow X_A && \text{(own-shared insufficiency)},\\
b_A(P_A)&\rightarrow X_A && \text{(own-private sufficiency)}.
\end{align}
Shared codes are additionally aligned by InfoNCE on $\mathrm{MLP}(S)$. Figure~\ref{fig:arch-m1} summarises the architecture.

\begin{figure}[htbp]
\centering
\includegraphics[width=0.98\textwidth]{M1.png}
\caption{Method~1 architecture. Per-modality encoders emit $(P_m,S_m)$. Predictors $h,f,g,k,b$ implement joint reconstruction, cross-shared sufficiency, private non-leakage, own-shared insufficiency, and private sufficiency. A contrastive head aligns projected shared codes.}
\label{fig:arch-m1}
\end{figure}

Let $\mathcal{L}_{\mathrm{std}}$, $\mathcal{L}_{\mathrm{shared}}$, $\mathcal{L}_{\mathrm{adv}}$, $\mathcal{L}_{\mathrm{own}}$, and $\mathcal{L}_{\mathrm{priv}}$ denote the sums over modalities of the corresponding reconstruction metrics. Training uses a two-step procedure per batch. Step~1 updates only the adversaries $g$ and $k$ on detached codes:
\begin{equation}
\mathcal{L}_{\mathrm{step1}}=\sum_A\bigl(\mathcal{L}_{\mathrm{adv}}^{A}+\mathcal{L}_{\mathrm{own}}^{A}\bigr).
\end{equation}
Step~2 updates encoders and constructive predictors with adversaries frozen:
\begin{equation}
\begin{aligned}
\mathcal{L}_{\mathrm{step2}}
&=
\sum_A\Bigl[
\mathcal{L}_{\mathrm{std}}^{A}
+\alpha\,\mathcal{L}_{\mathrm{shared}}^{A}
+\lambda_{\mathrm{priv}}\,\mathcal{L}_{\mathrm{priv}}^{A}
-\beta\,\mathcal{L}_{\mathrm{adv}}^{A}
-\gamma\,\mathcal{L}_{\mathrm{own}}^{A}
\Bigr]
+\lambda_{\mathrm{ctr}}\,\mathcal{L}_{\mathrm{ctr}}.
\end{aligned}
\end{equation}
The minus signs encourage encoders to make foreign private codes uninformative about~$X_A$ and to make own~$S_A$ insufficient for full reconstruction, while $b_A$ prevents the trivial collapse of~$P_A$ under privacy pressure. Checkpointing uses a constructive validation score that drops the signed adversarial terms, because raw $\mathcal{L}_{\mathrm{step2}}$ can become negative and unstable as an early-stopping signal. Across configurations, models are ranked by the label-free core validation loss $\mathcal{L}_{\mathrm{std}}+\mathcal{L}_{\mathrm{shared}}+\mathcal{L}_{\mathrm{priv}}+\mathcal{L}_{\mathrm{ctr}}$.

\subsection{Method 1 with VICReg ablation}

The incentive for this ablation is practical and diagnostic. Full Method~1 couples five predictors per modality with a two-step minimax procedure, and the own-shared adversary~$k$ (weight~$\gamma$) is particularly delicate: too strong a~$\gamma$ can strip useful shared content from~$S_A$, while too weak a~$\gamma$ leaves private bleed in~$S$. The Method~1+VICReg ablation therefore asks whether the \emph{cross-view} predictability story ($h$, $f$, $g$, and private sufficiency~$b$) already suffices when supplemented by a modern linear independence regulariser, without the brittle own-shared insufficiency term. Concretely, it removes~$k$ and adds a VICReg-style independence penalty between~$P_m$ and~$S_m$~\citep{bardes2022vicreg}. Predictors $h,f,g,b$ remain. The VICReg term combines variance hinges on both codes, within-code covariance whitening, and a cross-covariance penalty that targets linear $P\perp S$ independence. Figure~\ref{fig:arch-m1vic} illustrates the resulting architecture. Two-step training reduces to an adversary step on~$g$ only. Ranking again uses $\mathcal{L}_{\mathrm{std}}+\mathcal{L}_{\mathrm{shared}}+\mathcal{L}_{\mathrm{priv}}+\mathcal{L}_{\mathrm{ctr}}$, excluding $\mathcal{L}_{\mathrm{vic}}$ from the core rank so that independence regularisation is treated as a training regulariser rather than as the selection criterion. Empirically (Chapter~\ref{sec:results}), the ablation is more stable than full Method~1 and produces a cleaner subtype NMI ordering of~$S$ over~$P$, but it still fails the cross-view $\Delta\mathrm{X}\,R^{2}$ test---which motivates the simpler Method~2 and the spectral Method~3 that follow.

\begin{figure}[htbp]
\centering
\includegraphics[width=0.98\textwidth]{M1_VICReg.png}
\caption{Method~1+VICReg ablation. The own-shared adversary $k$ is removed; a VICReg $P\perp S$ term is added while retaining cross-view predictors $h,f,g$ and private sufficiency $b$.}
\label{fig:arch-m1vic}
\end{figure}

\subsection{Method 2: joint reconstruction and VICReg}

Method~2 abandons cross-reconstruction adversaries in favour of a simpler inductive bias. Each modality is reconstructed only from the concatenation of its own codes,
\begin{equation}
\hat{X}_m=D_m\bigl([P_m;S_m]\bigr),
\end{equation}
while VICReg encourages $P_m\perp S_m$ and InfoNCE aligns projected shared codes (Figure~\ref{fig:arch-m2}). The training objective is
\begin{equation}
\mathcal{L}
=
\lambda_{\mathrm{rec}}\,\mathcal{L}_{\mathrm{recon}}
+\lambda_{\mathrm{ctr}}\,\mathcal{L}_{\mathrm{ctr}}
+\lambda_{\mathrm{vic}}\,\mathcal{L}_{\mathrm{vic}},
\end{equation}
with all terms non-negative, so raw validation loss is safe for early stopping. Configurations are ranked by $\mathcal{L}_{\mathrm{recon}}+\mathcal{L}_{\mathrm{ctr}}$. Relative to MOCSS, Method~2 removes the requirement that~$S_m$ or~$P_m$ alone reconstruct~$X_m$, which is the most direct architectural fix of the dual-reconstruction failure mode.

\begin{figure}[htbp]
\centering
\includegraphics[width=0.98\textwidth]{M2.png}
\caption{Method~2 architecture. A single joint decoder reconstructs $X_m$ from $[P_m;S_m]$; VICReg replaces elementwise orthogonality; contrastive learning aligns shared codes.}
\label{fig:arch-m2}
\end{figure}

A companion variant, Method~2+OMIDIENT-private, replaces only the private encoder with OMIDIENT's product-Dirichlet branch while retaining Method~2's joint decoder, VICReg, and contrastive shared pathway. It is used to test whether Dirichlet private codes remain well-behaved once dual solo reconstruction is removed.

\subsection{Method 3: spectral shared targets with residual private decoding}

Method~3 supplies an explicit clustering-shaped target for shared codes. Offline, modality-wise $k$-nearest-neighbour graphs are fused and a normalised Laplacian embedding $U\in\mathbb{R}^{N\times r}$ with $r=16$ is computed~\citep{wang2014snf,ng2002spectral,vonluxburg2007tutorial}. A spectral head $f_{\mathrm{spec}}$ maps $\bar{S}=\mathrm{mean}_m(S_m)$ toward~$U$:
\begin{equation}
\mathcal{L}_{\mathrm{spec}}
=
\frac{1}{N}\sum_{i=1}^{N}
\bigl\|f_{\mathrm{spec}}(\bar{S}_i)-U_i\bigr\|_{2}^{2}.
\end{equation}
Private codes reconstruct residuals through
\begin{equation}
\hat{X}_m=\mathrm{Dec}_m\bigl(P_m,\mathrm{stopgrad}(S_m)\bigr),
\end{equation}
so gradients from residual decoding do not rewrite the shared code as a second full autoencoder of~$X_m$. Additional terms include a view-identity head on~$P$ (encouraging assay-specific content to remain in the private branch), optional VICReg, a gradient-reversal (GRL) leakage adversary that discourages private codes from predicting soft clusters of~$\bar{S}$, and a contrastive term whose weight $\lambda_{\mathrm{ctr}}$ is often selected as~$0$ under loss-based ranking. Figure~\ref{fig:arch-m3} shows the architecture. Core ranking uses $\mathcal{L}_{\mathrm{spec}}+\mathcal{L}_{\mathrm{res}}+\mathcal{L}_{\mathrm{ctr}}$. The unsupervised NMI of $k$-means on~$U$ itself, $\mathrm{NMI}(U)\approx 0.387$ on BRCA, is reported as a spectral ceiling. Full notation and formulae for the offline spectral target, spectral regression, view head, and GRL non-leakage block are collected in Appendix~\ref{app:method3}.

\begin{figure}[htbp]
\centering
\includegraphics[width=0.98\textwidth]{M3.png}
\caption{Method~3 architecture. An offline fused spectral embedding $U$ supervises mean shared codes via $f_{\mathrm{spec}}$; private decoders reconstruct residuals with stop-gradient on $S$; optional view-head and GRL leakage controls regularise role separation.}
\label{fig:arch-m3}
\end{figure}

\subsection{Training protocol}

All BRCA runs use batch size~$32$, up to~$100$ epochs, and early stopping with patience~$20$ on the method-specific validation score ($\mathrm{min\_delta}=0.005$). Learning rates and weight decays follow the MOCSS defaults used in the repository unless a sweep token indicates otherwise. Model selection is label-free: PAM50 labels never enter the ranking loss. After selection, clustering and probes are evaluated on the chosen configuration and aggregated over training seeds $\{21,22,23,24,25\}$.

\subsection{Evaluation protocol and interpretation of $\Delta\mathrm{X}\,R^{2}$}

The evaluation suite comprises:
\begin{enumerate}
\item subtype NMI, ARI, accuracy, and $k$NN accuracy on $\mathrm{mean}(S)$, $\mathrm{concat}(P)$, and $\mathrm{concat}(P,S)$;
\item cross-view $\Delta\mathrm{X}\,R^{2}$ and related Ridge probes;
\item own-view Ridge $R^{2}$ for $X\leftarrow S$, $X\leftarrow P$, and $X\leftarrow[P;S]$;
\item pairwise CKA among shared codes;
\item UMAP visualisations of $\mathrm{mean}(S)$ and $\mathrm{concat}(P)$.
\end{enumerate}

\paragraph{How to read $\Delta\mathrm{X}\,R^{2}$.}
A reliably positive gap means foreign shared codes explain another modality better than foreign private codes---evidence that cross-view signal concentrates in~$S$. Near-zero gaps indicate indistinct roles. Large positive gaps must still be checked against own-view $X\leftarrow P$: if that quantity is strongly negative, private codes have collapsed and the gap is not evidence of healthy privacy. Likewise, moderately high $\mathrm{concat}(P)$ NMI is acceptable when UMAP layouts differ from those of~$S$ and $\Delta\mathrm{X}\,R^{2}$ remains positive; private codes may then carry complementary subtype correlates rather than stolen shared geometry~\citep{levine2013integrated}.

\paragraph{Optional integration gap.}
An auxiliary quantity $\mathrm{gap}_{S-\mathrm{PS}}=\mathrm{mean}(S)\,\mathrm{NMI}-\mathrm{concat}(P,S)\,\mathrm{NMI}$ asks whether adding private codes hurts or helps the fused clustering embedding. A large positive gap means private content dilutes an already subtype-rich shared code; a negative gap means private codes contribute additional subtype structure to the concatenation. This thesis treats $\mathrm{gap}_{S-\mathrm{PS}}$ as secondary: it is useful context in compact tables, but it does not by itself certify role separation. The primary decomposition probes remain $\Delta\mathrm{X}\,R^{2}$, own-view $X\leftarrow P$, and the comparative UMAP geometries of~$S$ and~$P$.

\subsection{Strengths and limitations}

Table~\ref{tab:method-summary} summarises qualitative trade-offs. Method~1 offers the most explicit information-theoretic story but is heavy: fifteen predictors, two optimisers, and sensitive $(\alpha,\beta,\gamma,\lambda_{\mathrm{priv}})$ knobs. Method~1+VICReg reduces adversary complexity and adds linear independence, yet removing~$k$ weakens the insufficiency constraint on own~$S$. Method~2 is stable and simple, fixing MOCSS's dual-reconstruction flaw directly, but relies on VICReg and contrastive weights to prevent asymmetric collapse. Method~3 provides the strongest clustering-shaped shared target and the clearest empirical $\Delta\mathrm{X}\,R^{2}$ gains, at the cost of dependence on the quality of the offline fused graph and a spectral ceiling below the best concatenation NMI of reconstruction-based models.

\begin{table}[htbp]
\centering
\caption{Qualitative comparison of methods studied in this thesis (AE-private line unless noted).}
\label{tab:method-summary}
\begin{tabular}{p{2.4cm}p{3.2cm}p{3.2cm}p{3.6cm}}
\hline
 & Reconstruction & Independence / privacy & Main risk \\
\hline
MOCSS & dual solo $\mathrm{Dec}_S$, $\mathrm{Dec}_P$ & weak elementwise ortho & $P$ steals shared subtype \\
OMIDIENT & Dirichlet private + MOCSS shared & ortho optional; KL on $P$ & private-path collapse \\
Method~1 & $h,f,g,k,b$ predictability & adversaries + $b$ & minimax instability; weak concat NMI \\
M1+VICReg & $h,f,g,b$ + no $k$ & VICReg $P\perp S$ & low $\mathrm{concat}(P)$ $\neq$ disentangling \\
Method~2 & joint $D([P;S])$ & VICReg + contrastive & needs careful $\lambda$ balance \\
Method~3 & residual $\mathrm{Dec}(P,\mathrm{stopgrad}(S))$ & spectral $U$ + GRL leak & limited by $\mathrm{NMI}(U)$ ceiling \\
\hline
\end{tabular}
\end{table}
